\section{High-Level Language Design}
\label{sec:language}
\subsection{Dialects, Operations, and Rewrite Rules}
\label{sec:dialects}
A \emph{dialect} groups reusable operation definitions. Each \emph{operation}
declares its operands, result, attributes, and semantics. Attributes configure
an operation independently of tensor values.
Section~\ref{sec:attribute-roles} distinguishes the attribute roles.

A \emph{rewrite rule} imports operations from a dialect and describes a
replacement using \code{LHS => RHS}. Each side is a tensor expression written
as an s-expression. For example,
\code{(transpose (transpose X)) => (transpose X)} replaces two consecutive
transposes with one. The replacement's attributes determine whether the
rewrite preserves the result.

An operation occurrence can be named using \code{:}. For example,
\code{(d: dot X Y)} names the dot occurrence \code{d}. Later sections show
how \code{where} and \code{derive} use these names to read attributes,
express conditions, and construct replacement attributes.

\sidefigure[0.3]{figures/Axis, extent, and coordinate.pdf}{Axis, extent, and coordinate in a tensor of shape $(2,3,4)$.}{fig:attribute-roles}{A three-dimensional tensor with vertical axis 0 of extent 2, horizontal axis 1 of extent 3, and depth axis 2 of extent 4. The highlighted element has coordinate (0, 0, 2). Matching colors associate each axis with its shape and coordinate entries.}

\subsection{Axis, Extent, and Coordinate Attributes}
\label{sec:attribute-roles}

TEPL requires a type when declaring an operation attribute and distinguishes
three roles for integer attributes: axes, extents, and coordinates
(Figure~\ref{fig:attribute-roles}). These roles describe different properties
even when their integer values coincide. This distinction allows TEPL to
target only eligible axis fields for synthesis, while extent and coordinate
attributes remain supplied.
Appendix~\ref{app:language-declarations} gives declaration examples.

\begin{table}[!htbp]
\caption{Integer-attribute roles in TEPL.}
\label{tab:attribute-roles}
\small
\begin{tabularx}{\linewidth}{@{}llY@{}}\toprule
Role & Meaning & Example field\\\midrule
\code{axis} & Zero-based dimension position & Transpose \code{permutation: axis[]}\\
\code{extent} & Dimension size & Reshape \code{shape: extent[]}\\
\code{coordinate} & Position, offset, or stride & Slice \code{starts: coordinate[]}\\\bottomrule
\end{tabularx}
\end{table}

\sidefigure[0.5]{figures/Attribute Binding and Derivation.pdf}{Captured operation attributes can be read in \code{where} and \code{derive} blocks. The \code{derive} block explicitly defines RHS attribute transformation logic using these captured attributes.}{fig:attribute-binding}{An input graph with transpose and dot is replaced by a dot. Colored descriptor
bindings connect the rule to a table with two captured rows and one derived row.}

\subsection{Attribute Binding and Explicit Derivation}
\label{sec:binding}

LHS operation attributes are captured automatically at match time. Naming an operation
makes its descriptor accessible as \code{node.attrs}. For example,
\code{tr.attrs.permutation} selects the permutation of a named transpose
node \code{tr}.


Users can compute RHS attributes in a \code{derive} block by assigning a
descriptor to a named RHS node's \code{attrs} property, such as
\code{out.attrs}. Bracket notation provides syntactic sugar for directly
assigning an existing descriptor:

\mbox{\code{(dot[d.attrs] A B)}} reuses
\code{d.attrs} for the RHS dot. Descriptor compatibility is checked
separately from rewrite correctness.

For example, \code{gather} composes two captured transpose permutations:
\begin{lstlisting}[language=TEPL]
rule compose_transposes_explicit {
  (t2: transpose (t1: transpose X)) => (t3: transpose X)
  derive {
    t3.attrs = {
      permutation: gather(t1.attrs.permutation, t2.attrs.permutation)
    };
  }
}
\end{lstlisting}

\subsection{Conditions with \texorpdfstring{\code{where}}{where}}
\label{sec:where}


Some rewrite rules are applicable only under specific conditions. A
\code{where} block expresses these conditions as predicates combined with
\code{&&} and \code{||}.

For operations with declared semantics, acceptance establishes that the
attribute program is defined and produces a legal, equivalent replacement
whenever the condition holds (Section~\ref{sec:contract}). A condition can, for example, restrict a captured transpose permutation
to the identity. Appendix~\ref{app:condition-suggestions} gives the
batch-routing condition for generalized dot reassociation.

When verification fails, TEPL can suggest an additional sufficient condition
to add to \code{where}. The user adds the condition and checks the rule again
(Appendix~\ref{app:condition-suggestions}).

\subsection{Automatic Attribute Inference}
\label{sec:auto}
When eligible RHS axis fields are omitted and no derivation is provided, TEPL asks the
compiler to infer their transformation logic. The same transpose rewrite
can therefore be written without a \code{derive} block:
\begin{lstlisting}[language=TEPL]
rule compose_transposes_auto {
  (transpose (transpose X)) => (transpose X)
}
\end{lstlisting}
\sidefigure[0.5]{figures/Automatic Attribute Inference.pdf}{Omitted RHS axis fields trigger inference of a reusable attribute function at compile time. Runtime evaluates this function on the metadata captured for each match.}{fig:introducing-auto}{A vertical flow separates the user rule and operation semantics, compiler
synthesis and verification, and runtime application of the generated function.}
The compiler infers the explicit permutation transformation above once;
runtime evaluates it on each match's captured metadata
(Figure~\ref{fig:introducing-auto}). Unsupported inference produces a
diagnostic. Section~\ref{sec:problem} formalizes the contract across ranks
and shapes.
\FloatBarrier
